\section{Introduction}

\subsection{Motivation and Contributions}

\HE{} studies recovery from public homotopy-class data when the Hodge operator
whose kernel is the message space is not part of the public input.  In a
known-lattice closest-vector problem, the lattice is presented before the
search begins.  In HCP, the recovery target is
\[
  h=P_H(b)\in\ker(\Delta_1),
\]
but the public class does not include the degree-one Hodge--Laplacian
\(\Delta_1\), its kernel, or the restriction maps from which \(\Delta_1\) is
built.  The private
recovery object is therefore not a scalar secret or a vector in a published
lattice; it is the harmonic cochain in the kernel of a hidden Hodge operator.

The operator dependence can be summarized as follows. A hidden representative
determines the restriction data, which determine the adjacent
coboundaries and hence the degree-one Hodge operator:
\[
  x
  \longmapsto
  \{R_{\sigma\to\tau}^{x}\}_{\sigma\prec\tau}
  \longmapsto
  (\delta_{0,x},\delta_{1,x})
  \longmapsto
  \Delta_{1,x}
  \longmapsto
  \bigl(\ker\Delta_{1,x},P_{H_x}\bigr).
\]
The recovery operation is Hodge-theoretic; ``homotopical'' describes the
framework in which multiple hidden representatives may occupy the same public
class while inducing different operators. The distinction is important: it is
hidden-operator ambiguity, not whether a displayed matrix happens to have
binary entries, that drives HCP.

This changes the comparison problem. A real linear span can be realized as the
kernel of a suitable linear operator, but a discrete lattice cannot equal a
nonzero real kernel. If an explicit lattice \(\Lambda\) is supplied inside a
harmonic subspace \(W=\ker(\Delta_W)\), orthogonal projection of the target
onto \(W\) preserves the closest vectors of \(\Lambda\). That bridge does not
make public-data HCP a known-lattice problem: public homotopy-class data
may be compatible with distinct hidden operators, distinct kernels, and
different harmonicity judgments.  The paper isolates this boundary and then
separates it from the finite verifier-facing projection artifacts used later.
The public classifier supplies the realized class value \(\alpha\).
The full cochain \(b=h+e\) is separate from the accepted exact commitment
payload: its decoded exact cochain has zero harmonic projection in its
specified Hodge system. A bound-observation relation ties the full sample to that payload,
its digest, snapshot and lifecycle context. Releasing \((\alpha,b)\)
requires an explicit disclosure or a specified extraction of that full
sample; a proof about the exact payload alone does not expose it.
The selected observation need not present \(\Delta_1\) or
\(\ker(\Delta_1)\). If full Pl\"ucker coordinates of the operative kernel
are additionally disclosed, however, they determine its fixed-metric
projector and hence the HCP target from \(b\).
The fibration ID is the spectral fingerprint
\(\operatorname{Spec}_{\mathrm{multi}}(\Delta_1)\), obtained only after the
relevant Hodge data have been recovered or certified; it is not part of the
public classifier value.

The principal notation is fixed globally as follows. Symbols used only inside
a single proposition or construction are defined locally.
\begin{table}[ht]
  \centering
  \small
  \begin{tabular}{@{}p{0.27\linewidth}p{0.65\linewidth}@{}}
    \hline
    Symbol & Meaning \\
    \hline
    \(p:E\to B\) & Helfer-style fibration projection \\
    \(\mathcal P_X(x,y)\) & path object from \(x\) to \(y\) in \(X\) \\
    \(\pi:A\to P\) & path-respecting public classifier \\
    \(\alpha=\pi(x)\in P\) & realized public classifier value \\
    \(A_\alpha=[x]_\pi\) & fiber of hidden representatives with public label \(\alpha\) \\
    \((\alpha,b)\) & public HCP observation \\
    \(\Delta_A:A\to A\times A\) & diagonal map in the fibrational semantics \\
    \(P_X\), \(F:P_X\to\mathsf{Vect}_{\R}\) & finite cell poset and its face-to-coface cellular sheaf \\
    \(a\in C^0,\ c=\delta_0a\in C^1\) & degree-zero opening and committed exact one-cochain \\
    \(\delta_0,\delta_1,\Delta_1\) & adjacent coboundaries and degree-one Hodge--Laplacian \\
    \(\ker(\Delta_1),\ P_H\) & harmonic subspace and harmonic projection \\
    \(\Delta_W=I-P_W\) & kernel-normalized projection operator associated with \(W\) \\
    \(\mathcal L,\ \ell\) & local syndrome target and affine MCA line, respectively \\
    \hline
  \end{tabular}
  \caption{Core notation used across the semantic and HCP layers.}
  \label{tab:core-notation}
\end{table}

This paper develops this boundary through the following statements.
\begin{enumerate}[label=(\roman*)]
  \item The Harmonic Cochain Problem (HCP): recovery of the operative
    component \(h=P_Hb\) for a hidden degree-one Hodge operator \(\Delta_1\).
  \item A Hodge-structured noise model with cochain decomposition
    \(b = h + e\), where \(h \in \ker(\Delta_1)\) and
    \(e \in \image(\delta_0) \oplus \image(\delta_1^\ast)\).
  \item A CVP comparison separating field-linear image from discrete lattice:
    generated subspaces can be realized as kernels, and, when
    \(\Lambda\subseteq W\) is explicitly supplied, projecting the target onto
    \(W\) preserves the discrete closest-vector condition on \(\Lambda\).
  \item A hidden-operator ambiguity theorem: publicly equivalent Hodge operators
    may disagree on harmonicity, so the CVP comparison does not make a correct
    explicit-lattice CVP oracle into a uniform public-data HCP solver.
  \item Constructed integral examples in the finite signed-incidence model:
    one homological public observation has distinct nonzero harmonic targets;
    a weighted pair shares \(b=(1,2)\) with integer coefficients two and one,
    with a unique CVP answer for each privately presented lattice.
  \item The supporting formal framework:
    fibrational identity, homotopy classes, exact Hodge/cochain invariants,
    spectral/Tsirelson admissibility, and Heyting-style public judgments.
  \item A projection chain through Pl\"ucker coordinates and cochain commitments to a
    canonical Reed--Solomon endpoint, with a conditional far-separation gate and
    an MCA bad-challenge bound imported at an explicit cited theorem boundary.
\end{enumerate}

Finite Hodge theory follows the cellular-cochain framework of
Hansen--Ghrist \cite{hansenGhrist2018sheaves}; the algebraic model used here is
specified in \cref{sec:hcp-definition}. Kernel realization and the
Pythagorean projection identity are elementary linear-algebraic foundations.
Helfer's homotopy-invariant semantics \cite{helfer2019fol} motivates the public
classifier; the constructed examples prove only the explicitly named
nonzero-cohomology-test fragment. The selected formalization integrates these
objects with the constructed observation ambiguities and recovery boundaries.
It establishes neither an asymptotic CVP-hardness reduction nor a sampled
encryption-security theorem. The independently stated Lean entry covers the
six Hodge/CVP groups and two observation-boundary groups in its source
crosswalk; the broader spectral, oracle and coding interfaces retain their
separately stated assumptions and formalization status.

The companion manuscript \emph{Capacity-Radius Proximity for Smooth-Domain
Reed--Solomon Codes} \cite{vanDerVlugt2026proximity} develops the universal
list-decoding and MCA results independently of HCP. The present paper supplies
the HCP foundations and the application-specific projection interface. Shared
coding definitions connect the two papers; the companion's coding theorems do
not supply the source-to-code far-separation witness required here.

The corresponding Lean definitions, reductions, gates, theorem boundaries,
paper-to-Lean cross-reference, and file-by-file source coverage are recorded in
the appendices.
